\documentclass[instructions.tex]{subfiles}
\begin{document}
 
\chapter{Qui sont ceux qui vivent long-temps.}
 

\primo On ne devrait compter le nombre de ses jours, que selon qu'ils sont bien employés. La gloire de la vieillesse et du grand âge n'est pas le nombre des années, mais la vie sainte; parce qu'on ne doit estimer avoir vécu qu'autant qu'on s'est rendu digne de vivre, et qu'il vaudrait mieux n'être pas que de vivre mal. Une vie pure et innocente est donc la seule qui mérite d'être appelée longue\footnote{Aetas senectutis, vita immaculata. Sap.4}. 

Celui qui est enlevé de ce monde en sa jeunesse, avant que la malice ait corrompu son esprit, n'a pas laissé que de remplir la course d'une longue vie\footnote{Consumatus in brevi, explevit tempora multa. Sap.13.}. Parce qu'il a rempli ses jours, il a donné les fruits de sainteté qu'il pouvait donner. Au contraire, celui qui, étant avancé en âge, a vécu dans le désordre, ne doit pas compter avoir vécu, parce que ce n'est pas vivre, que de vivre sans servir Dieu. Un arbre n'est pas vivant, du moins il est comme mort, lorsqu'il ne donne ni ombrage, ni fruit à son maitre. Le roi Saül régna long-temps; cependant l'Écriture ne compte que deux années de son règne, parce qu'il ne régna pas toujours selon les desseins de Dieu. 

Sur ce principe, pouvez-vous dire que jusqu'ici vous ayez vécu ? Vous passez trente ou quarantes ans, et peut-être n'avez-vous pas encore vécu une année. On vous croyait vivant, mais vous n'aviez qu'une apparence de vie; vous étiez mort devant Dieu, puisque vous ne produisiez aucune action de vie pour le salut\footnote{Nomen habes quod vivas, et mortuus es. Apoc.3.}. 

Un saint religieux, âgé de soixante-dix ans, interrogé combien d'années il avait vécu en religion, répondit:\og Je doute fort si j'ai vécu un seul jour\fg. On ne demandera pas au jour du jugement combien vous aurez vécu d'années, mais comment vous les aurez passées\footnote{Non quam diu, sed quam benè. Imit.}. Toutes les années que vous avez vécu hors de la grâce de Dieu, dans le péché, sont autant d'années perdues, dont on ne tiendra compte que pour vous punir. La vie ne vous est donnée que pour le salut; c'est donc comme si vous n'aviez pas vécu, si vous l'avez employée à d'autres choses. 

\secundo N'est-il pas temps de commencer à vivre et de réparer vos pertes ? Quoiqu'on ne puisse rappeler le temps, on peut néanmoins, en quelque sorte, réparer la perte qu'on en a faite. Un voyageur qui s'est arrêté sur la route, double le pas pour arriver à son terme. Un artisan travaille avec plus d'activité, pour regagner les momens qu'il a perdus. Faites de même; redoublez le pas vers le ciel, travaillez avec ardeur pour y arriver. Vos maux sont grands, mais ils ne sont pas incurables. 

Si vous avez beaucoup perdu en perdant le passé, vous pouvez beaucoup gagner pour le ciel, en profitant du présent. Servez-vous des grâces que Dieu vous donne; avec votre coopération, elles peuvent vous faire, de grand pécheur, un grand saint. Mais le temps presse; si vous différez, vous risquez tout. Il vous reste peu de temps à vivre; préparez-vous à ce passage terrible qui décide de tout et pour toujours. Donnons au moins à Dieu, dit Saint Pierre Chrysologue, le reste d'une vie que nous avons employée jusqu'ici aux choses du siècle; nous avons donné le passé aux commodités du corps, donnons au moins quelques jours au salut de l'âme\footnote{Vivamus deo paululum, qui seculo viximusl totum dedimus corpori annum, demus animae paucos dies.}. 

La résolution est prise: je veux commencer à vivre. Le moins que je doive, après tant de temps perdu, c'est d'employer le peu de temps qui me reste à contenter Dieu et à me procurer le salut. Il faut que je me sauve; je n'ai que cette seule affaire au monde: si elle réussit, tout est gagné pour moi. Je viens donc, ô mon Dieu vous offrir les misérables restes d'une vie jusqu'à présent si mal employée. Quoique j'aie mal vécu, j'ai encore cette confiance en votre miséricorde, que vous ne rebuterez point un coeur contrit et humilié. Voilà les sentimens d'une âme qui désire réparer le temps perdu.


\section{Résolutions}

J'ai peu vécu jusqu'ici, parce que mes années ont été remplies de péchés, et presque vides de bonnes oeuvres. Il faut que je répare le passé:

\primo En évitant avec soin les péchés auxquels je suis si sujet. Quels sont-ils ... ?

\secundo En faisant le bien qui m'est commandé, mais en le faisant dans des vues saintes, non plus par vanité, orgueil, vain désir d'être estimé, loué, applaudi des hommes. 

\prayer{Mon Dieu, faites que je commence enfin à vivre pour vous et pour le ciel.}

\end{document}
